\section{Charged derivatives and the finite Hecke action}
\label{sec:charged-theta-interface}

The neutral principal line in AC-G35 does not make its tangent
representation neutral.  The relevant question is whether a
nontrivial charge in the differentiated theta space changes
the finite character that enters the source cohomology bound.
The supplied source answers one precisely scoped version of
this question.

\begin{proposition}[AC-G36: exact derivative charge and a finite-character obstruction]
\label{AC-G36}
Use the source harmonic-polynomial space
$\operatorname{Harm.rowHighest}(c,c')$.
Write the row degrees, in its literal notation, as
\[
 d_{\tau,a}=\operatorname{partOf}(c_\tau)(a),\qquad
 d'_{\tau,b}=\operatorname{partOf}(c'_\tau)(b).
\]
Choose $r_\tau\in\{1,-1\}$ with
$\tau(\omega)=\omega^{r_\tau}$.
Scaling each differential direction by the action of
$[\omega]$ scales every polynomial differential operator
from this highest-weight space by
\[
 \omega^{q(c,c')},\qquad
 q(c,c')=\sum_\tau r_\tau
       \left(\sum_a d_{\tau,a}-\sum_b d'_{\tau,b}\right)
                        \quad\text{in }\mathbb Z/3.   \tag{G36.1}
\]
The relative cubic charge is independent of the common
metaplectic vacuum phase; in the source Fock normalization
that phase can be a sixth root of unity.

Fix the source datum $D$, its two Hermitian spaces, finite
Weil representation and finite Schwartz/level data.
Changing only the creation polynomial $P$ preserves every
finite Hecke eigensystem on a nonzero created vector.
In particular it cannot change a fixed one-dimensional
finite character $\mu$ to $\chi_{269}\mu$ when
$\chi_{269}$ is nontrivial.
This assertion does not compare different ranks, finite
Schwartz functions, levels or Arthur packets.
\end{proposition}
\begin{proof}
\proofstep{1. The literal highest-polynomial charge.}
The exact definitions and support theorem are
\texttt{Harm.rowHighest} and
\path{Harm.mem_rowHighest_iff_support} in
\path{Definitions/ThetaValues/ExistsMemRowHighestTotalDegreeEqTwo.lean}.
Every monomial $X^m$ in a polynomial $P$ in that space has
\[
 \sum_jm(x_{\tau,a,j})=d_{\tau,a},\qquad
 \sum_jm(y_{\tau,b,j})=d'_{\tau,b}.                    \tag{G36.2}
\]
The file
\path{Definitions/ThetaValues/PolyDerivSmulLetter.lean}
identifies the positive-row type with $\tau$ and the
negative-row type with its conjugate, in the exact theorems
\texttt{rowType\_rowOf\_inl} and
\texttt{rowType\_rowOf\_inr}.
Thus $[\omega]$ multiplies the associated $x$ directions
by $\omega^{r_\tau}$ and the $y$ directions by
$\omega^{-r_\tau}$.
By (G36.2), the product of these scalars for a monomial is
$\omega^{q(c,c')}$, independent of that monomial.

For clarity, this last step is also a direct proof, without
an analytic transformation assumption.
An iterated derivative is multilinear in its differential
directions.  Scaling those directions by scalars $s_j$
multiplies the term for $X^m$ by $\prod_js_j^{m_j}$.
When this product is the same scalar for every term in $P$,
it factors out of the sum defining $\operatorname{polyDeriv}$.
This is precisely the supplied theorem
\path{GlobalTheory.ThFr.polyDeriv_smul_letter}
(the namespace is that of the named source file).
It proves (G36.1).
The equation is a statement about scaling directions;
a convention using pullback on cotangent coordinates
inverts this character, consistently on both sides.

The polynomial derivative is part of the actual source
construction.  In \path{SpanningFamilies/IsSystem.lean},
\texttt{GlobalTheory.SW.thP} is \texttt{polyDeriv} applied
to the theta function \texttt{thA} and evaluated at zero.
In \path{SpanningFamilies/IsStabForm.lean}, the fields
\texttt{Dic.dir\_eq}, \texttt{Dic.mem} and \texttt{Dic.span}
identify the algebraic directions and highest-polynomial
space with that analytic construction.
Those definitions do not themselves alter a finite
Schwartz function or introduce a finite Galois character.

\proofstep{2. Keep the metaplectic half-weights.}
The exact Fock formulas in \texttt{DataTheta.lean} are
\texttt{omE\_diag}, \texttt{cartan1} and \texttt{cartan2}.
For a mode $j$ of sign $s_j\in\{1,-1\}$,
they give the diagonal infinitesimal weight
$s_j(m_j+1/2)$ on $X^m$.
For a diagonal compact path with integer rotation rates
$a_j$, integration therefore gives, at angle $2\pi/3$,
\[
 X^m\longmapsto
 \underbrace{\exp\left(\frac{\pi i}{3}\sum_j s_ja_j\right)}_{\displaystyle\beta\in\mu_6}
 \omega^{\sum_j s_ja_jm_j}X^m.                        \tag{G36.3}
\]
Here integration is justified directly by the source
\texttt{FockData.der} and \texttt{GrpData.expt\_add}:
the group law moves the derivative at zero to every point
of the one-parameter group.  On an eigenmonomial the
resulting differential equation is
$v'(t)=i\lambda v(t)$ with
$\lambda=\sum_js_ja_j(m_j+1/2)$.
Differentiating $e^{-i\lambda t}v(t)$ gives zero, which
proves the displayed exponential action.
This uses no choice of an unproved global lift.

The same path acts on the vacuum by $\beta$.
Dividing its monomial eigenvalue by the vacuum eigenvalue
leaves exactly the integral cubic weight in (G36.3).
For example a single positive mode of rate one has
$\beta=e^{\pi i/3}$, whose cube is $-1$; dropping this
half-weight would incorrectly replace that chosen
order-six lift by an order-three action.
Conversely, this common phase does not change the relative
polynomial charge and does not specify a character of
finite Frobenius elements.
Formula (G36.3) is an archimedean representation calculation;
it is not an identification of its vacuum with every
algebraic rigidified principal theta line of AC-G35.

\proofstep{3. Creation commutes with the actual finite action.}
For fixed $P$, abbreviate the source creation operator by
\[
 C_Pf=\operatorname{pureTensorO}(\ldots,f,P).
\]
The full supplied theorem
\path{GlobalTheory.AdForm.FinComm.finite_comm_creation}
in \path{Proofs/DecoratedLatticeProofs/FiniteCommCreation.lean}
states, with its standard-pair and Weil-data hypotheses,
\[
 \Omega(g,h)C_Pf=C_P\Omega(g,h)f,
       \qquad g\in G_f(S),\ h\in G_f(T).              \tag{G36.4}
\]
The more general theorem
\path{GlobalTheory.OrientedHeis.om_fin_pureTensorO}
in \path{Proofs/TwistedOperators/CompAtMtOfMemGfC.lean}
has the same conclusion with the displayed scale and frame
parameters retained.
Its preceding theorem \texttt{compAt\_mt\_of\_mem\_GfC}
explains the mechanism: the finite subgroup is the kernel
of the archimedean projection, so its component at each
archimedean place is the identity.
The full proof of (G36.4) consequently commutes that finite
action with each infinitesimal creation operator; induction
commutes it with every monomial and then every $P$.
This is a source theorem about the actual action, rather
than an analogy with an abstract oscillator.

A finite Hecke operator on vectors fixed by a compact open
subgroup is a finite linear combination of the relevant
finite-group translates.  By linearity, (G36.4) implies
\[
                       T C_P f=C_P T f.              \tag{G36.5}
\]
If $Tf=a_Tf$, then $T(C_Pf)=a_T(C_Pf)$.
Thus every nonzero $C_Pf$ has exactly the old finite
Hecke eigenvalues.  The same argument gives
$C_P(T-a_T)^n=(T-a_T)^nC_P$, so it also preserves the
specified generalized eigenspaces; it need not be injective.

For a one-dimensional unramified character, the eigenvalue
at Frobenius is nonzero.  A nontrivial cubic character is
nontrivial on Frobenius at some place outside any prescribed
finite excluded set, by the finite Galois correspondence
and Chebotarev applied to its cyclic splitting field.
At that place the values of $\mu$ and $\chi_{269}\mu$ differ.
Equation (G36.5) therefore prevents a nonzero change from
one to the other by $C_P$ alone.
This argument concerns a one-dimensional character;
it makes no assertion that an arbitrary higher-rank
representation cannot have a self-twist.

\proofstep{4. The supplied jet identity and the remaining gate.}
The theorem
\path{GlobalTheory.KJInf.CoreG.thetaKernel_eq_jet_thC}
in \path{ThetaKernelSplitPlacesEqJetClassicalTheta.lean}
uses (G36.4) explicitly.  Its conclusion identifies the
created theta kernel with a scalar times
$\operatorname{jetP}(\operatorname{thC}(\Phi,\ldots),P)$,
with the same finite Schwartz data $\Phi$ and the same $P$.
Therefore replacing the kernel by that literal jet
expression does not evade the commutation calculation.
In particular a nonzero value of (G36.1) alone cannot
replace a scalar block by a cubic finite-character block
in \texttt{NData}.

To obtain such a block one must also modify and compare
finite data or an actual descent action, and then verify
the new Hecke, local inertia and cohomological identities.
A change of ranks or of packet falls outside (G36.4), which
compares a single fixed source datum; no impossibility
claim for such a new comparison is made.
Together with the exact \texttt{chi\_def} cancellation and
\texttt{sha\_eq} restriction of AC-G29, this identifies the
remaining source-specific task rather than supplying the
missing \texttt{LinkN} theorem.
\end{proof}
